\section{Discussion}

The results support an audit-and-deployment story rather than a universal
repair story. Ordinary discrimination can substantially overstate the
trustworthiness of answer-visible process verifiers; answer-matched
counterfactual evaluation exposes this gap; and the most reliable deployment
intervention in the present setting is to remove direct answer access.

\paragraph{Established findings.}
Three conclusions are stable across the empirical package. First,
\amcd{} is materially more informative than ordinary AUROC for understanding
whether a verifier is useful in downstream selection. Second, judge-style
evaluators are consistently answer-sensitive and therefore should not be
interpreted as process-grounded simply because they discriminate valid from
invalid traces on average. Third, stronger model scale does not make
answer-visible verification safe by default: the strongest deployment model in
the study is a masked reranker, and the strongest external-source verifier is
currently a step-only frozen baseline rather than an answer-visible reranker.

\paragraph{What remains unsupported.}
The results do not show that all trained verifiers are dominated by outcome
leakage, nor do they show that a simple repair term uniformly beats a strong
masked baseline. \texttt{visible\_cond\_swap} improves answer invariance and
can move \amcd{} on harder slices, but it does not dominate masked baselines or
masked reranking. The dual-head line is likewise informative but negative at
the current scale. The right read is therefore not ``repair failed,'' but
``repair sharpened the mechanism story without becoming the deployment answer.''
In this paper, repair terms are best read as diagnostic interventions rather
than as a new architecture family.

\paragraph{Why both benchmark regimes matter.}
\textsc{Craft-Deploy} and \textsc{Craft-Clean} answer different questions. The
first is the better deployment stress test: it exposes selection gain,
exploitability, and attacker-transfer differences sharply. The second is the
better faithfulness test: it preserves the visible-vs-masked process gap while
reducing the benchmark artifacts that would otherwise blur interpretation. The
divergence between these regimes is therefore a result, not a nuisance or a
post-hoc benchmark swap.

\paragraph{Open issues.}
Two limits remain important. External human blind audit on \textsc{Craft-Clean}
is still pending, so that benchmark should be described as strongly audited
rather than fully closed. By contrast, \textsc{ProcessBench} is now an
operational external-source package that is closed for its intended
corroboration role rather than an exploratory side result. The remaining
benchmark-closure gap is therefore specifically about \textsc{Craft-Clean}
blind audit, not about whether the story survives on real step-level data. On
the deployment side, planning-style naturalized traces remain the hardest
regime for the best current model, and even after fully materializing a small
external answer-matched local-repair subset, the visible view still does not
beat the masked one. The external counterfactual evidence therefore supports
the same conservative claim boundary. It does not support a new visible-method
recovery story.
